\subsection{EXTENSION BY CONTINUITY; DOUBLE LIMIT}

THEOREM 1. Let X be a topological space, A a dense subset of X, \(f: \mathbf{A} \to \mathbf{Y}\) a mapping of A into a regular space Y. A necessary and sufficient condition for \(f\) to extend to a continuous mapping \(\overline{f}: \mathbf{X} \to \mathbf{Y}\) is that, for each \(x \in \mathbf{X}\) , \(f(y)\) tends to a limit in Y when \(y\) tends to \(x\) while remaining in A. The continuous extension \(\overline{f}\) of \(f\) to X is then unique.

The uniqueness of \(\vec{f}\) follows from the principle of extension of identities (no. 1, Proposition 2, Corollary 1). It is clear that the condition is necessary, for if \(\vec{f}\) is continuous on X, then for each \(x \in X\) we have

\[
\overline {{f}} (x) = \lim _ {y > x, y \in \mathbf {A}} \overline {{f}} (y) = \lim _ {y > x, y \in \mathbf {A}} f (y)
\]

(§ 7, no. 5). Conversely, suppose that the condition is satisfied and define

\[
\overline {{f}} (x) = \lim _ {y \rightarrow x, y \in \mathbf {A}} f (y)
\]

for each \(x \in X\) ; \(\overline{f}(x)\) is a well-defined element of Y, since Y is Hausdorff. We have to show that \(\overline{f}\) is continuous at each point \(x \in X\) . Let then \(V'\) be a closed neighbourhood of \(\overline{f}(x)\) in Y; then by hypothesis there is an open neighbourhood V of x in X such that \(f(V \cap A) \subset V'\) . Since V is a neighbourhood of each of its points, we have

\[
\overline {{f}} (z) = \lim _ {y > z, y \in \mathrm{V} \cap \mathrm{A}} f (y)
\]

for each \(z \in \mathbf{V}\) , and from this it follows that \(\overline{f}(z) \in f(\mathbf{V} \cap \mathbf{A}) \subset \mathbf{V}'\) , since \(\mathbf{V}'\) is closed. The result now follows from the fact that the closed neighbourhoods of \(f(x)\) form a fundamental system of neighbourhoods of \(f(x)\) in \(\mathbf{Y}\) .

The mapping \(\bar{f}\) is said to be obtained by extending f by continuity to X.

In the statement of Theorem 1 the hypothesis that Y is regular cannot be weakened without imposing additional restrictions on X, A or f (Exercise 19).

COROLLARY. Let \(\mathfrak{F}_1\) be a filter on a set \(X_1\) , and \(\mathfrak{F}_2\) a filter on a set \(X_2\) ; let \(\mathfrak{F}_1 \times \mathfrak{F}_2\) be the product filter (§ 6, no. 7) on \(X = X_1 \times X_2\) , and let \(f\) be a mapping of \(X\) into a regular space \(Y\) . Suppose that

\begin{enumerate}
\def\labelenumi{\alph{enumi})}
\item
  \(\lim_{\mathfrak{F}_1\times \mathfrak{F}_2}f\) exists.
\item
  \(\lim_{x_2, \mathfrak{F}_2} f(x_1, x_2) = g(x_1)\) exists for all \(x_1 \in \mathbf{X}_1\) .
\end{enumerate}

Then \(\lim_{x_1, \mathfrak{F}_1} g(x_1)\) exists and is equal to \(\lim_{\mathfrak{F}_1 \times \mathfrak{F}_2} f\) .

Let \(X_1' = X_1 \cup \{ \omega_1 \}\) (resp. \(X_2' = X_2 \cup \{ \omega_2 \}\) ) be the topological space associated with the filter \(\mathfrak{F}_1\) (resp. \(\mathfrak{F}_2\) ) (§ 6, no. 5, Example). In the product space \(X' = X_1' \times X_2'\) let \(X''\) be the union of the subspaces \(X_1 \times X_2'\) and \(\{ (\omega_1, \omega_2) \}\) . \(X\) is clearly a dense subspace of \(X''\) , and the hypotheses imply that \(f(y_1, y_2)\) tends to a limit when \((y_1, y_2)\) tend to any point \((x_1, x_2)\) of \(X''\) whilst remaining in \(X\) . The existence of the extension of \(f\) by continuity to \(X''\) then follows from Theorem 1. Since also \((\omega_1, \omega_2)\) lies in the closure of \(X_1 \times \{ \omega_2 \}\) relative to \(X''\) , the result follows immediately (§ 7, no. 5).

